\chapter{Marco Conclusivo}
\label{ch:marco_conclusivo}

El presente capítulo resume los resultados obtenidos a lo largo del desarrollo
de la plataforma de simulación neuromórfica, formula recomendaciones orientadas
a futuras iteraciones del proyecto y delinea las líneas de trabajo futuro que
se derivan de las limitaciones identificadas.

\section{Conclusiones}
\label{sec:conclusiones_principales}

El trabajo desarrolló \texttt{neurolab}, un simulador modular de
memristores, neuronas LIF, plasticidad y crossbar, y lo conectó con el
agente Robot Sim. La evidencia se limita a simulaciones, verificaciones
internas, soluciones analíticas y curvas digitalizadas atribuidas a
publicaciones. No constituye una validación general de hardware físico.
El estado de cumplimiento de los objetivos, con el alcance de cada
evidencia, se resume en la Tabla~\ref{tab:cumplimiento_objetivos}.

\begin{enumerate}
    \item \textbf{Modelo memristivo y variabilidad:} El modelo de Strukov
    alcanza $R^2=0.999657$ para $x(t)$ y $R^2=0.996994$ para $I(t)$ bajo la
    alineación y escalado especificados en el Capítulo~3; para resistencia
    y conductancia se obtienen $R^2=0.998319$ y $R^2=0.983726$,
    respectivamente. D2D/C2C están habilitados en las corridas actuales de
    Robot Sim mediante modificadores por celda. Las pruebas estadísticas
    reproducen propiedades de distribuciones sintéticas, no variabilidad
    medida en dispositivos.
    \item \textbf{Histéresis y estado:} Se implementaron ciclos I--V y
    evolución acotada de estado. Para Prezioso, se presenta una
    reconstrucción visual aproximada basada en la Fig.~1(b). El artículo
    disponible no especifica completamente el barrido temporal y la
    digitalización no conserva tiempo ni etiquetas de rama; no se calculan
    métricas de ajuste ni se demuestra la validación I--V por rama. El
    objetivo se considera parcialmente cumplido.
    \item \textbf{Neurona LIF:} La implementación integra, dispara, reinicia y
    respeta el período refractario. La curva $f-I$ se contrastó con la solución
    analítica del mismo modelo ($R^2=0.999964$, MAE $=0.086$ Hz,
    RMSE $=0.151$ Hz), no con datos experimentales independientes.
    \item \textbf{Plasticidad:} Se implementó una regla STDP en el módulo de
    sinapsis; al ajustar nueve puntos y evaluar cuatro reservados, la ventana
    obtiene $R^2=0.94220$ y RMSE $=2.136$ puntos porcentuales. El holdout es
    pequeño. La curva digitalizada de Jo asigna signo positivo a LTP, pero
    la definición temporal en prosa del artículo usa el signo contrario;
    el proyecto declara y prueba la convención
    $\Delta t=t_{\mathrm{post}}-t_{\mathrm{pre}}$. Para LTP/LTD se emplean
    los pulsos publicados de $+3.2$ y $-2.8\,\mathrm{V}$, de
    $300\,\mu\mathrm{s}$, y lectura a $1\,\mathrm{V}$. En los 18 puntos de
    test se obtuvo $R^2=0.73988$, Pearson $r=0.99303$,
    RMSE $=8.618$ nA y ERP $=31.23\%$;
    las métricas por rama son deficientes ($R^2=-8.54$ para LTP y
    $R^2=-829.99$ para LTD). Los parámetros efectivos se ajustan a la
    misma curva digitalizada, por lo que no son una caracterización
    independiente del dispositivo. La evidencia es parcial. Las
    corridas de Robot Sim resumidas en este documento usan REINFORCE con
    baseline. El software entregado ofrece tambi\'en R-STDP en los cuatro runners y sus
pruebas funcionales verifican partes de su mecanismo. La carpeta de evidencias
incluye un piloto exploratorio de 28 corridas en CartPole, visi\'on y obst\'aculos
(dos semillas y dos valores de $X_0$), adem\'as de cuatro corridas R-STDP de Car
bajo un protocolo inicial. Una repetici\'on separada de Car R-STDP con
evaluaci\'on cada 100 episodios us\'o semillas 2 y 4. La muestra sigue siendo
peque\~na, Car REINFORCE no est\'a emparejado en estos lotes y el protocolo Car
inicial difiere; no se afirma superioridad estad\'istica ni se generaliza a
otras condiciones.
    \item \textbf{Crossbar e integración:} El solver MNA incluye resistencias
    de línea y sensado en lectura. En el caso $2\times2$ la diferencia
    relativa de norma respecto a la VMM ideal es $10.98\%$; es una
    comprobación numérica, no experimental. Robot Sim usa una matriz
    abstracta $X$ cuyos valores determinan $R(X)$ y $G(X)$. En ambos modos, el memristor vol\'atil de salida modula la corriente
    que llega a la LIF. En REINFORCE, la acci\'on se decodifica desde las
    corrientes de columna; en el R-STDP actual tambi\'en se decodifica desde
    esa pol\'itica softmax y WTA no la restringe. La elegibilidad R-STDP usa
    eventos pre del codificador y la acci\'on decodificada centrada por la
    softmax; los spikes LIF se registran, pero no act\'uan como evento post
    en el piloto cuantitativo del runner. Este piloto usa elegibilidad
    pre/acci\'on centrada por la pol\'itica y modulada por retornos descontados;
    no es la regla can\'onica de pares pre/post. Se distinguen tres resultados:
    la STDP local del laboratorio usa spikes pre/post y trazas temporales;
    una variante de tarea con spikes LIF reales como evento post obtuvo bajo
    rendimiento exploratorio; y las corridas de la carpeta de evidencias,
    etiquetadas R-STDP por el modo del software, corresponden al h\'ibrido
    pre/acci\'on. Estas corridas no demuestran que la R-STDP can\'onica
    aprenda las tareas. Las reglas cambian en software, sin alterar la
    topolog\'ia ni el modelo el\'ectrico del crossbar. El bajo desempe\~no de
    la variante directa est\'a observado, pero las pruebas no aíslan una causa
    \'unica.
    \item \textbf{Validación cuantitativa:} Las métricas se calcularon con
    procedimientos de alineación y limitaciones documentados sobre curvas
    digitalizadas atribuidas a publicaciones. La evidencia es heterogénea:
    Prezioso cuenta con una reconstrucción visual aproximada sin métricas
    de ajuste; STDP se evalúa en cuatro puntos reservados y LTP/LTD muestra
    errores elevados, en particular por rama. Los parámetros LTP/LTD se
    estiman sobre la misma curva digitalizada. Ninguna comparación
    representa validación frente a mediciones experimentales propias.
\end{enumerate}

\section{Respuesta a la pregunta de investigación}
\label{sec:respuesta_pregunta}

La concordancia depende de la magnitud y del protocolo. Con
$R^2=1-\mathrm{SSE}/\mathrm{SST}$, el modelo local de Strukov muestra alta
concordancia de forma con las series digitalizadas usadas ($R^2=0.999657$
para el estado, $0.996994$ para la corriente, $0.998319$ para la
resistencia y $0.983726$ para la conductancia). Los ERP de conductancia
Strukov ($19.04\%$) y de la curva STDP ($14.87\%$) son errores máximos
normalizados por el rango de sus respectivas series; el umbral ERP
$<10\%$ del Capítulo~1 se aplica a parámetros puntuales, mientras la
aceptación de curvas se evalúa con $R^2$ y coincidencia cualitativa.
Por ello, ambos ERP no contradicen por sí solos el criterio de curva:
Strukov cumple $R^2>0.90$, y STDP obtiene $R^2=0.94220$ en un holdout
pequeño. Para Prezioso se dispone de una reconstrucción visual aproximada
de la Fig.~1(b), no de métricas de ajuste: el barrido temporal no está
completamente especificado y los datos digitalizados no permiten separar
las ramas SET/RESET. La histéresis se mantiene como evidencia parcial.
En STDP, el holdout de cuatro puntos tiene RMSE $=2.136$ puntos
porcentuales; su tamaño no permite una afirmación amplia de
generalización. Además, la curva digitalizada asigna $\Delta t>0$ a LTP,
mientras que la definición temporal en prosa de Jo et al. presenta el
signo opuesto. El proyecto usa explícitamente
$\Delta t=t_{\mathrm{post}}-t_{\mathrm{pre}}$ y prueba que $\Delta t>0$
potencia y $\Delta t<0$ deprime.
En LTP/LTD, se obtuvo en 18 puntos de test
$R^2_{\mathrm{residual}}=0.73988$, MAE $=8.349$ nA,
RMSE $=8.618$ nA y ERP $=31.23\%$, aplicando los pulsos y la lectura
indicados en la Fig.~2a. Pearson del test agregado es $r=0.99303$
($r^2=0.98610$), pero este coeficiente mide covariación y no sustituye
al error residual: la predicción de LTP es constante (Pearson no
definido) y LTD obtiene $r=-0.22170$. Las métricas desagregadas son débiles
($R^2=-8.54$ para LTP y $R^2=-829.99$ para LTD); el resultado agregado
no basta para demostrar concordancia por rama. Los parámetros efectivos
se ajustaron a la misma curva digitalizada, por lo que no acreditan
amplitud ni predicción independientes. En LIF, Euler
concuerda con la solución analítica del propio modelo
($R^2=0.999964$). En conjunto, la evidencia cuantitativa es
heterogénea: alta en Strukov y LIF, limitada en STDP por el tamaño del
holdout y no satisfactoria por rama en LTP/LTD. La
variabilidad D2D/C2C se probó estadísticamente con muestras
sintéticas y se habilitó en el crossbar del agente, sin comparación con
medidas de hardware. El MNA se comprobó en un caso pequeño, cuya
diferencia relativa de norma respecto a la VMM ideal fue $10.98\%$; no
es una comparación con hardware. En consecuencia, la evidencia
disponible no permite afirmar concordancia experimental general para
dispositivos, redes SNN o crossbars físicos.
Entre las limitaciones de integración están el uso de $X$ como estado
del modelo simulado y no como dispositivo físico, la LIF como monitor/WTA
en REINFORCE y la ausencia de programación con \textit{sneak paths}. Los
resultados principales de Auto 2D de esta tesis usan una semilla; la carpeta
de evidencias agrega un piloto independiente con tres semillas emparejadas
por condición y evaluación final congelada. El piloto comparativo posterior ya cubre CartPole, visi\'on y obst\'aculos,
pero solo con dos semillas; Car cuenta con una repetici\'on R-STDP separada y
no incluye una condici\'on REINFORCE emparejada. Falta ampliar la muestra y
uniformar el protocolo Car antes de estimar diferencias entre reglas. En el
agente, $X_0=0.1$ (equivalente a $x_0=0.9$) es una elecci\'on de dise\~no que
difiere de la convenci\'on de estado inicial del dispositivo aislado; el
piloto tambi\'en explora $X_0=0.9$, pero no sustituye un estudio de sensibilidad
con mayor tama\~no muestral ni permite extrapolar a otras inicializaciones.

Otras limitaciones del alcance son: las referencias de literatura se
introducen como curvas digitalizadas y no como mediciones propias; el
holdout STDP comprende solo cuatro puntos y los parámetros efectivos de
LTP/LTD se estiman sobre la misma curva de referencia, con desempeño
deficiente por rama; la reconstrucción Prezioso es visual y no identifica
el protocolo temporal del artículo; en REINFORCE, la LIF actúa como
monitor/WTA y no como decisor principal;
la matriz $X$ representa el estado de celdas simuladas, no un crossbar
    físico; las actualizaciones del agente cambian directamente el estado
    abstracto $X$, con límites y variabilidad configurada, sin simular una
    forma de onda de programación ni pulsos discretos por actualización; los
    valores de $\gamma$ y
la actualización de elegibilidad difieren entre tareas; el solver MNA
incluye resistencias de línea y el acoplamiento de la topología de
lectura simulada, pero no la configuración clásica de
\textit{sneak paths} con líneas flotantes ni la programación parásita;
y los resultados principales de Auto 2D cuentan con una sola semilla. El
piloto separado de la carpeta de evidencias amplía esa tarea a tres semillas
por condición, pero sigue siendo exploratorio. Por tanto, los resultados de
las tareas no constituyen una comparación general entre algoritmos ni
permiten generalizar el rendimiento más allá de las configuraciones
evaluadas.

\section{Recomendaciones}

A partir de las lecciones aprendidas durante el desarrollo y de las limitaciones
observadas en la validación, se formulan las siguientes recomendaciones:

\begin{itemize}
    \item \textbf{No idealidades de crossbar:} la lectura incluye
    resistencias de línea y el acoplamiento de la topología simulada, pero no
    la configuración clásica de \textit{sneak paths} con líneas flotantes ni
    la escritura parásita. Se recomienda modelar y validar ambos casos antes
    de extrapolar a arreglos de mayor densidad.
    \item \textbf{Acoplamiento térmico:} complementar el ruido estocástico
    macroscópico actual con un modelo de calentamiento de Joule dependiente de
    la disipación temporal a escala nanométrica.
    \item \textbf{Robustez del solver:} el esquema de integración Euler
    Forward requiere pasos temporales restrictivos ($dt \leq 10^{-3}$\,s) en
    los umbrales ultra no lineales del dispositivo; se recomienda evaluar
    integradores adaptativos (p.\ ej.\ Runge-Kutta) para ampliar la ventana
    numérica de simulación.
    \item \textbf{Reproducibilidad continua:} mantener la suite de validación
    como herramienta de referencia en \texttt{neuromorphic\_lab/validaciones/},
    ejecutando automáticamente los scripts (p.\ ej.\
    \texttt{python 01\_strukov\_2008.py}) bajo Python 3.10+ para revalidar el
    ecosistema en cada modificación, sin intervención manual del operador.
    \item \textbf{Validación de escritura y selección:} contrastar con datos
    experimentales los pulsos de escritura y los cambios de $X$ y $G$ del
    modo R-STDP; evaluar si la LIF debe participar en la selección de
    acciones del modo REINFORCE. Mantener explícita la distinción entre
    $X$, conductancia y pesos normalizados.
    \\item \\textbf{Cobertura estad\'istica:} ampliar el piloto REINFORCE/R-STDP
    desde dos semillas a una muestra previamente dimensionada; emparejar Car
    bajo el mismo intervalo de evaluaci\'on y reportar resultados por condici\'on.
    Mantener separados los resultados incompletos y las corridas que alcanzan
    el presupuesto m\'aximo sin {\'e}xito. Documentar adem\'as las fuentes y el
    proceso de digitalizaci\'on de cada CSV atribuido a la literatura.
\end{itemize}

\section{Trabajo Futuro}

El simulador se evaluó en cuatro tareas: control con recompensa densa
(CartPole-v1), clasificación supervisada (Visión $8\times8$),
navegación con recompensa escasa (Gridworld $10\times10$) y control
en un entorno 2D de estado continuo (Auto 2D). La extensión a otros
entornos de control continuo y a tareas con acciones continuas se
plantea como trabajo futuro. También se consideran la unificación de
la regla de elegibilidad entre tareas, la integración de la LIF en el
lazo de decisión y la modelización de \textit{sneak paths} durante
escritura.

Las líneas de desarrollo identificadas sobre la base de los resultados
obtenidos comprenden:

\begin{itemize}
    \item \textbf{Crossbars de alta densidad:} extender el solver MNA para
    arreglos de escala industrial (por ejemplo $128 \times 128$) con
    selectores (\textit{memristor switches}) y modelado completo de sneak
    paths y margen de ruido.
    \item \textbf{Redes neuronales pulsantes completas:} combinar la
    neurona LIF implementada con las sinapsis STDP en redes SNN capaces de
    tareas de reconocimiento de patrones de mayor complejidad (p.\ ej.\
    MNIST o sensores DVS).
    \item \textbf{Comparación de reglas de aprendizaje:} completar y
    documentar una comparación controlada de REINFORCE y R-STDP en las
    cuatro tareas de Robot Sim, con semillas emparejadas, presupuesto de
    entrenamiento equivalente, checkpoints seleccionados solo con
    validación y pruebas finales con pesos congelados. Publicar los
    registros originales y las métricas por semilla antes de concluir
    que una regla supera a la otra.
    \item \textbf{Co-simulación con hardware:} exportar las matrices
    conductivas entrenadas a flujos de diseño de circuitos (p.\ ej.\
    SPICE o frameworks de computación en memoria) para acercar la simulación
    a la implementación física.
    \item \textbf{Benchmarking inter-simulador:} comparar las métricas de
    Neuromorphic Lab con herramientas de referencia del área para
    contextualizar su precisión dentro del estado del arte.
\end{itemize}
